\documentclass[11pt,a4paper,reqno]{amsart}

\usepackage[utf8]{inputenc}
\usepackage[T1]{fontenc}
\usepackage{amsmath,amssymb,amsfonts,amsthm,mathtools}
\usepackage{mathrsfs}
\usepackage{microtype}
\usepackage{booktabs}
\usepackage[dvipsnames]{xcolor}
\usepackage{hyperref}
\usepackage{geometry}
\usepackage{enumitem}
\usepackage{cite}
\usepackage{graphicx}

\geometry{
  top=2.7cm,
  bottom=2.7cm,
  left=2.7cm,
  right=2.7cm,
  headheight=14pt
}

\hypersetup{
  colorlinks=true,
  linkcolor=NavyBlue,
  citecolor=Mahogany,
  urlcolor=TealBlue,
  pdfauthor={Reinaldo M. Silva-Filho},
  pdftitle={Emergent Spacetime, Holographic Entanglement, and Quantum Geometry: Proved Results, Conjectures, and Observational Limits}
}

% --- Theorem Environments ---
\newtheorem{theorem}{Theorem}[section]
\newtheorem{lemma}[theorem]{Lemma}
\newtheorem{proposition}[theorem]{Proposition}
\newtheorem{corollary}[theorem]{Corollary}

\theoremstyle{definition}
\newtheorem{definition}[theorem]{Definition}
\newtheorem{example}[theorem]{Example}
\newtheorem{remark}[theorem]{Remark}
\newtheorem{axiom}[theorem]{Axiom}
\newtheorem{conjecture}[theorem]{Conjecture}

\numberwithin{equation}{section}

% --- Custom Notations ---
\providecommand{\R}{\mathbb{R}}
\newcommand{\C}{\mathbb{C}}
\newcommand{\N}{\mathbb{N}}
\newcommand{\Z}{\mathbb{Z}}
\providecommand{\Sph}{\mathbb{S}}
\providecommand{\doi}[1]{\href{https://doi.org/#1}{\nolinkurl{doi:#1}}}
\providecommand{\Hyp}{\mathbb{H}}
\newcommand{\Id}{\mathrm{Id}}
\newcommand{\Tr}{\operatorname{Tr}}
\newcommand{\rank}{\operatorname{rank}}
\newcommand{\supp}{\operatorname{supp}}
\newcommand{\Lip}{\operatorname{Lip}}
\providecommand{\BV}{\operatorname{BV}}
\providecommand{\Area}{\operatorname{Area}}
\newcommand{\Ric}{\operatorname{Ric}}
\providecommand{\II}{\mathrm{I\!I}}
\providecommand{\cutnorm}[1]{\|#1\|_{\square}}
\newcommand{\norm}[1]{\left\|#1\right\|}
\newcommand{\inner}[2]{\left\langle #1, #2 \right\rangle}
\newcommand{\dif}{\,\mathrm{d}}
\newcommand{\abs}[1]{\left|#1\right|}

\title[Emergent Spacetime and Quantum Geometry]{Emergent Spacetime, Holographic Entanglement,\\ and Quantum Geometry:\\ Proved Results, Conjectures, and Observational Limits}

\author{Reinaldo Maia Silva-Filho}
\address{PPGEE/DES, Universidade Federal de Lavras (UFLA) \\
Lavras, Minas Gerais, Brazil}
\email{reinaldo.filho1@estudante.ufla.br}
\thanks{This research was financed in part by the Coordena\c{c}\~ao de Aperfei\c{c}oamento de Pessoal de N\'ivel Superior -- Brasil (CAPES) -- Finance Code 001.}
\date{\today}


\DeclareMathOperator*{\argmin}{argmin}
\providecommand{\doi}[1]{\href{https://doi.org/#1}{\nolinkurl{doi:#1}}}
\providecommand{\Hyp}{\mathbb{H}}
\providecommand{\Sph}{\mathbb{S}}
\providecommand{\II}{\mathrm{II}}
\providecommand{\R}{\mathbb{R}}
\providecommand{\Area}{\mathrm{Area}}
\providecommand{\Hsn}{\mathcal{H}}
\providecommand{\BV}{\mathrm{BV}}
\providecommand{\TV}{\mathrm{TV}}
\providecommand{\cutnorm}[1]{\|#1\|_{\square}}

\begin{document}

\begin{abstract}
This chapter collects the points of contact between three programs developed in the preceding chapters: (i) continuous multilinear tensor networks and information geometry (Part I), (ii) fractional calculus on continuous Pascal simplices (Part II), and (iii) $L^\infty$-minimax extrinsic curvature of submanifolds (Part III), and examines what each of them can and cannot say about problems of quantum gravity. We separate proved statements from conjectures throughout.

The proved results are elementary but exact:
(1) for a heat kernel with Fourier symbol $|k|^{2\alpha}$ on $\R^m$ the spectral dimension is $d_s = m/\alpha$, and for the assumed isotropic two-scale symbol $k^2 + \ell_P^2 k^4$ in four dimensions we give a closed-form expression for $d_s(\tau)$ interpolating between $4$ and $2$, the same end values as the flow found in Causal Dynamical Triangulations and in Ho\v{r}ava--Lifshitz gravity (where it arises from a different, anisotropic operator);
(2) on a spacelike hypersurface whose second fundamental form is bounded in operator norm by $\kappa^*$, the vacuum Hamiltonian constraint confines the scalar curvature to the sharp interval $2\Lambda - 6(\kappa^*)^2 \le {}^{(3)}R \le 2\Lambda + 2(\kappa^*)^2$;
(3) the second-order Taylor coefficient of the continuous multinomial entropy at the barycenter is the Euclidean metric of the $A_{m-1}$ root space, written in the basis of simple roots through the Cartan matrix.
We also restate, with its standard hypotheses, the linearized derivation of the Einstein equations from the first law of entanglement, which is imported from the literature.
The remaining links --- restriction of loop quantum gravity states to Jordan loops, mean curvature flow as a search procedure for Ryu--Takayanagi surfaces, condensation of graphon Ricci flows to smooth manifolds, and the relation between Kac--Rice complexity and horizon entropy --- are formulated as explicit conjectures. A phenomenological section records geometric parametrizations of Standard Model quantities together with their free-parameter content, and the observational section gives order-of-magnitude estimates showing that Planck-suppressed signatures of the kind discussed here lie far below the reach of current and planned experiments.
\end{abstract}

\maketitle

\tableofcontents
\newpage

% =========================================================================
\section{Introduction and the Tripartite Unification Framework}
\label{sec:intro}
% =========================================================================

Theoretical physics currently rests upon two extraordinarily successful yet mutually incompatible pillars:
\begin{enumerate}[label=(\alph*)]
    \item \textbf{General Relativity (GR):} A geometric field theory of classical pseudo-Riemannian spacetimes $(M^4, g_{\mu\nu})$ where gravity is manifest as curvature, governed by the non-linear Einstein field equations $G_{\mu\nu} + \Lambda g_{\mu\nu} = 8\pi G_N T_{\mu\nu}$.
    \item \textbf{Quantum Field Theory (QFT):} A linear, operator-algebraic framework defined on infinite-dimensional Hilbert spaces $\mathscr{H}$, requiring fixed background spacetimes and canonical commutation relations.
\end{enumerate}

When these two frameworks are confronted at the Planck scale $\ell_P = \sqrt{\frac{\hbar G}{c^3}} \sim 1.616 \times 10^{-35}\text{ m}$, severe mathematical breakdowns occur: perturbative non-renormalizability, the Wheeler--DeWitt "problem of time", quantum foam topology fluctuations, and the black hole information scrambling paradox.

This chapter does not resolve these problems. Its aim is narrower: to record precisely which statements about them follow from the mathematics of the preceding chapters, and which remain conjectural. Three programs are involved:
\begin{itemize}
    \item \textbf{Program I: Multilinear Tensor Geometry and Information Holography} \cite{silvafilho2026functional, silvafilho2026flows, silvafilho2026spacetime}: the quantum Fisher information metric of continuous tensor networks (cMPS/cMERA) and its relation to emergent bulk geometry, following the entanglement-thermodynamics program of \cite{faulkner2014gravitation, faulkner2017nonlinear}.
    \item \textbf{Program II: Simplicial Fractional Calculus on Pascal Simplices} \cite{silvafilho2026pascal}: continuous multinomials via Euler's $\Gamma(z)$ and fractional simplicial operators, used here as models of scale-dependent diffusion of the kind that underlies the running spectral dimension observed in Causal Dynamical Triangulations (CDT) \cite{ambjorn2005spectral, carlip2017dimension}.
    \item \textbf{Program III: Minimax Extrinsic Curvature} \cite{silvafilho2026noneuclidean}: $L^\infty$ bounds on the second fundamental form of spacelike hypersurfaces and their consequences for the ADM constraints.
\end{itemize}
Statements labelled \emph{Theorem} or \emph{Proposition} are proved here or in the cited chapters; statements labelled \emph{Conjecture} are not.

% =========================================================================
\section{Simplicial Foundations and the Fractional Spectral Dimension of Quantum Spacetime}
\label{sec:simplicial_fractional}
% =========================================================================

In discrete quantum gravity models—including Regge Calculus, Spin Foams, and Causal Dynamical Triangulations (CDT) \cite{ambjorn2005spectral}—spacetime is approximated by simplicial complexes composed of $m$-simplices $\Delta_m$.

\subsection{Analytic Continuation of Pascal's Simplex}
Following \cite{silvafilho2026pascal}, let the continuous standard $(m-1)$-dimensional simplex of scale $x > 0$ be:
\begin{equation}
    \label{eq:standard_simplex}
    \Delta_{m-1}(x) \coloneqq \left\{ \mathbf{t} = (t_1, \dots, t_m) \in \R_{\ge 0}^m \colon \sum_{i=1}^m t_i = x \right\}.
\end{equation}
The continuous multinomial density over $\Delta_{m-1}(x)$ is defined via Euler's Gamma function:
\begin{equation}
    \label{eq:continuous_multinomial}
    \binom{x}{\mathbf{t}} \coloneqq \frac{\Gamma(x + 1)}{\prod_{i=1}^m \Gamma(t_i + 1)}.
\end{equation}
The total volume integral $I_m(x) \coloneqq \int_{\Delta_{m-1}(x)} \binom{x}{\mathbf{t}} \dif \sigma(\mathbf{t})$, where, as in Chapter~3, $\dif\sigma = \dif t_1 \cdots \dif t_{m-1}$ is Lebesgue measure on the projection of $\Delta_{m-1}(x)$ to the first $m-1$ coordinates (with the $(m-1)$-dimensional Hausdorff measure the integral is larger by the factor $\sqrt{m}$), grows like $m^x$: with the normalized volume $\mathcal{J}_m(x) \coloneqq m^{-x} I_m(x)$, Theorem~3.2 of Chapter~3 proves $\lim_{x \to \infty} \mathcal{J}_m(x) = 1$ by Laplace's method, so $I_m(x) = m^x\,(1 + o(1))$ \cite{silvafilho2026pascal}. For $m = 2$ Theorem~2.1 of Chapter~3 also gives a trigonometric integral for $\mathcal{J}_2$; for $m \ge 3$ Chapter~3 shows that the analogous torus integral does not represent $I_m$, and no oscillatory representation of $\mathcal{J}_m$ is used.

\subsection{The Fractional Simplicial Laplacian and Running Spectral Dimension}
For $\alpha > 0$ let $(-\Delta)^\alpha$ denote the Fourier multiplier with symbol $|k|^{2\alpha}$ on $\R^m$. For $\alpha \in (0,1)$ it has the singular-integral representation
\begin{equation}
    \label{eq:frac_simplicial_laplacian}
    (-\Delta)^\alpha f(\mathbf{t}) = C_{m,\alpha}\, \mathrm{P.V.}\!\int_{\R^m} \frac{f(\mathbf{t}) - f(\mathbf{s})}{\|\mathbf{t} - \mathbf{s}\|^{m + 2\alpha}} \dif \mathbf{s}, \qquad C_{m,\alpha} = \frac{4^\alpha\, \Gamma\!\left(\frac{m}{2}+\alpha\right)}{\pi^{m/2}\, |\Gamma(-\alpha)|},
\end{equation}
while for $\alpha \ge 1$ only the multiplier definition is used (for $\alpha = 1, 2$ it is the local operator $-\Delta$, respectively $\Delta^2$). The simplicial Beta-kernel operators of Chapters~3--4 \cite{silvafilho2026pascal} are \emph{not} of this form. Their kernels are integrable with compact support, so their symbols are bounded (Chapter~3, Theorem~7.2), and they do not produce a power-law return probability at short times. The results below concern the model symbol $|k|^{2\alpha}$ and do not follow from the simplicial operators.

Numerical simulations of Causal Dynamical Triangulations \cite{ambjorn2005spectral} show a spectral dimension decreasing from $d_s \approx 4$ at large diffusion times to $d_s \approx 2$ at short ones; a similar short-distance value appears in several other approaches \cite{carlip2017dimension}. In Ho\v{r}ava--Lifshitz gravity with anisotropic scaling exponent $z$ in $3+1$ dimensions the ultraviolet value is $d_s = 1 + 3/z$, so $d_s = 2$ requires $z = 3$ \cite{horava2009spectral}. The following elementary computation records how the spectral dimension depends on the order of the symbol.

\begin{proposition}[Spectral Dimension of a Homogeneous Symbol]
\label{thm:running_spectral_dimension}
Let $P(\tau; \mathbf{x}, \mathbf{y})$ be the heat kernel of $\partial_\tau u + (-\Delta)^\alpha u = 0$ on $\R^m$, $\alpha > 0$, and define $d_s \coloneqq -2\, \dif \log P(\tau;\mathbf{x},\mathbf{x}) / \dif \log \tau$. Then
\begin{equation}
    \label{eq:spectral_dimension_formula}
    d_s(\alpha, m) = \frac{m}{\alpha}
\end{equation}
for every $\tau > 0$. In particular $d_s(1,4) = 4$ and $d_s(2,4) = 2$.
\end{proposition}

A homogeneous symbol therefore gives a \emph{constant} spectral dimension; a running spectral dimension requires a symbol with two scales, treated in closed form in Section~\ref{sec:analytical_solutions}. That the value $d_s = 2$ at short scales is associated with power-counting renormalizability is a heuristic argument used in Ho\v{r}ava--Lifshitz gravity \cite{horava2009spectral}; it is not established here. The continuum value is also distinct from the harmonic Sierpi\'nski value $d_s = \frac{2\ln(m+1)}{\ln(m+3)} \approx 1.654$ ($m=4$) of Chapter~6, which combines the renormalization factor $m+3$ derived there (Proposition~2.2 of Chapter~6) with the Kigami--Lapidus theorem.

For $\alpha > 1$ the heat kernel takes negative values, and $P(\tau; \mathbf{x}, \mathbf{x})$ is then the diagonal of the kernel rather than a return probability; the definition of $d_s$ applies unchanged.

\begin{proof}
The symbol is $\lambda_k = |k|^{2\alpha}$. The diagonal of the heat kernel (the return probability when $\alpha \le 1$) after diffusion time $\tau$ is:
\begin{equation}
    P(\tau; \mathbf{x}, \mathbf{x}) = \frac{1}{(2\pi)^m} \int_{\R^m} e^{-\tau k^{2\alpha}} \dif^m k = \frac{\operatorname{Vol}(\Sph^{m-1})}{(2\pi)^m} \int_0^\infty k^{m-1} e^{-\tau k^{2\alpha}} \dif k.
\end{equation}
Substituting $u = \tau k^{2\alpha}$, whence $k = (u/\tau)^{\frac{1}{2\alpha}}$ and $\dif k = \frac{1}{2\alpha \tau} (u/\tau)^{\frac{1}{2\alpha} - 1} \dif u$:
\begin{equation}
    P(\tau; \mathbf{x}, \mathbf{x}) = \frac{\operatorname{Vol}(\Sph^{m-1})}{2\alpha (2\pi)^m} \tau^{-\frac{m}{2\alpha}} \int_0^\infty u^{\frac{m}{2\alpha} - 1} e^{-u} \dif u = C(m, \alpha) \tau^{-\frac{m}{2\alpha}}.
\end{equation}
Hence $\log P = \log C(m,\alpha) - \frac{m}{2\alpha}\log\tau$, and at every $\tau > 0$
\begin{equation}
    d_s = -2\, \frac{\dif \log P(\tau; \mathbf{x}, \mathbf{x})}{\dif \log \tau} = -2 \left( -\frac{m}{2\alpha} \right) = \frac{m}{\alpha},
\end{equation}
which establishes \eqref{eq:spectral_dimension_formula}.
\end{proof}

\begin{proposition}[Barycentric Hessian and the $A_{m-1}$ Root Metric]
\label{thm:cartan_metric_emergence}
Let $\mathbf{t}^* = (x/m, \dots, x/m)$ be the barycenter of $\Delta_{m-1}(x)$ and $\mathbf{v} \in T_{\mathbf{t}^*}\Delta_{m-1} = \{\mathbf{v} \in \R^m : \sum_i v_i = 0\}$. Then
\begin{equation}
    \label{eq:cartan_quadratic}
    \begin{gathered}
    \log \binom{x}{\mathbf{t}^* + \mathbf{v}} = \log \binom{x}{\mathbf{t}^*} - Q(\mathbf{v}) + \mathcal{O}(|\mathbf{v}|^3), \\
    Q(\mathbf{v}) = \frac{\psi'(x/m + 1)}{2} \sum_{i=1}^m v_i^2 = \frac{m}{2x}\left(1 + \mathcal{O}(x^{-1})\right) \sum_{i=1}^m v_i^2 .
    \end{gathered}
\end{equation}
Writing $\mathbf{v} = \sum_{j=1}^{m-1} c_j \boldsymbol{\alpha}_j$ in the simple roots $\boldsymbol{\alpha}_j = \mathbf{e}_j - \mathbf{e}_{j+1}$ of $A_{m-1}$, one has $\sum_i v_i^2 = \mathbf{c}^T \mathbf{A}_{m-1} \mathbf{c}$, where $\mathbf{A}_{m-1}$ is the Cartan matrix of $A_{m-1} \cong \mathfrak{sl}(m)$, so that $Q = \frac{m}{2x}(1+\mathcal{O}(x^{-1}))\, \mathbf{c}^T \mathbf{A}_{m-1} \mathbf{c}$.
\end{proposition}
\begin{proof}
The linear term vanishes on $T_{\mathbf{t}^*}\Delta_{m-1}$ because all partial derivatives $-\psi(t_i+1)$ coincide at the barycenter. The Hessian of $-\sum_i \log\Gamma(t_i+1)$ is $-\operatorname{diag}(\psi'(t_i+1))$, which equals $-\psi'(x/m+1)\,\Id$ at $\mathbf{t}^*$, and $\psi'(z+1) = z^{-1} + \mathcal{O}(z^{-2})$. Finally, $\mathbf{A}_{m-1}$ is the Gram matrix $(\boldsymbol{\alpha}_i \cdot \boldsymbol{\alpha}_j)$ of the simple roots, which gives $\sum_i v_i^2 = |\mathbf{v}|^2 = \mathbf{c}^T \mathbf{A}_{m-1} \mathbf{c}$.
\end{proof}

\begin{remark}
The content of Proposition~\ref{thm:cartan_metric_emergence} is that the Hessian at the barycenter is a multiple of the Euclidean metric on the sum-zero hyperplane. Any $S_m$-invariant quadratic form on that hyperplane has the same property, since $S_m$ acts irreducibly on it, and the hyperplane is the Cartan subalgebra of $\mathfrak{sl}(m)$ with the Weyl group $S_m$. The appearance of the Cartan matrix is therefore a statement about the permutation symmetry of the multinomial. It does not by itself produce gauge fields or gauge dynamics.
\end{remark}

% =========================================================================
\section{Minimax Foliations, Gravitational Shear, and the Hamiltonian Constraint}
\label{sec:minimax_adm}
% =========================================================================

In the canonical Arnowitt--Deser--Misner (ADM) $3+1$ formalism of General Relativity, the spacetime metric is decomposed as:
\begin{equation}
    \label{eq:adm_metric}
    \dif s^2 = -N^2 \dif t^2 + \gamma_{ij}(\dif x^i + N^i \dif t)(\dif x^j + N^j \dif t),
\end{equation}
where $N$ is the lapse function, $N^i$ is the shift vector, and $\gamma_{ij}$ is the induced Riemannian metric on Cauchy slices $\Sigma_t$.

The conjugate momentum is $\pi^{ij} = \frac{\sqrt{\gamma}}{16\pi G_N}(K^{ij} - K \gamma^{ij})$, where $K_{ij} = \frac{1}{2N}(\partial_t \gamma_{ij} - \nabla_i N_j - \nabla_j N_i) = \frac12 \mathcal{L}_n \gamma_{ij}$ is the extrinsic curvature (the second fundamental form $\II_{\Sigma_t}$), $\nabla_i$ is the Levi-Civita connection of $\gamma$, and $n$ is the future unit normal. With this sign, which is the one for which $\pi^{ij}$ takes the form above, expanding slices have $K > 0$. Chapter~8 \cite{silvafilho2026noneuclidean} uses the opposite sign, $K_{ij} = -\frac12\mathcal{L}_n\gamma_{ij}$, so the two tensors differ by a factor $-1$. The statements about Cauchy slices below involve $K_{ij}$ only through $|K_{ij}v^iv^j|$, $K_{ij}K^{ij}$ or $K^2$, and are therefore the same in both conventions. The subsection on thin shells concerns timelike hypersurfaces and uses Israel's convention $K_{ab} = h_a{}^\mu h_b{}^\nu \nabla_\mu n_\nu$, as in the section on junction conditions of Chapter~8.

The classical Hamiltonian constraint (whose quantization is the Wheeler--DeWitt equation) reads $\mathcal{H}_{\mathrm{ADM}} = \frac{16\pi G_N}{\sqrt{\gamma}}\left( \pi_{ij}\pi^{ij} - \frac{1}{2}\pi^2 \right) - \frac{\sqrt{\gamma}}{16\pi G_N}\left( {}^{(3)}R - 2\Lambda \right) + \sqrt{\gamma}\,\rho = 0$, with $\rho$ the energy density measured by the normal observer. In terms of the extrinsic curvature it takes the geometric form
\begin{equation}
    \label{eq:wdw_hamiltonian}
    {}^{(3)}R - 2\Lambda - 16\pi G_N \rho = K_{ij} K^{ij} - K^2 = \sigma_{ij}\sigma^{ij} - \frac{2}{3}K^2,
\end{equation}
where $\sigma_{ij} \coloneqq K_{ij} - \frac{1}{3} K \gamma_{ij}$ is the trace-free part of $K_{ij}$. We write $\rho = 0$ for vacuum.

\subsection{Minimax Extrinsic Curvature and Pointwise Constraint Bounds}
In numerical relativity, poorly chosen slicings can drive $K_{ij}$ to large values and produce coordinate pathologies, which motivates slicing conditions that control the extrinsic curvature.

Following the non-Euclidean variational theory of \cite{silvafilho2026noneuclidean}, we formulate the $L^\infty$ \emph{Minimax Extrinsic Curvature Problem}:
\begin{equation}
    \label{eq:minimax_hypersurface_problem}
    \kappa^*(\mathcal{M}, g) \coloneqq \inf_{\Sigma \in \operatorname{Cauchy}(\mathcal{M})} \|\II_\Sigma\|_{L^\infty(\Sigma, g)} = \inf_\Sigma \sup_{x \in \Sigma} \sup_{v \in T_x \Sigma, |v|=1} |K_{ij}(x) v^i v^j|.
\end{equation}

\begin{theorem}[Pointwise Constraint Bounds on Slices of Bounded Extrinsic Curvature]
\label{thm:minimax_hamiltonian_regularization}
Let $(\mathcal{M}^4, g_{\mu\nu})$ be a globally hyperbolic spacetime satisfying the Einstein equations, and let $\Sigma^* \subset \mathcal{M}$ be a spacelike hypersurface with $\|\II_{\Sigma^*}\|_{L^\infty} \le \kappa^* < \infty$ (for instance, a minimizer of \eqref{eq:minimax_hypersurface_problem}, \emph{if} one exists with finite value). Then, at every $x \in \Sigma^*$:
\begin{enumerate}[label=(\roman*)]
    \item the quadratic invariant of the extrinsic curvature satisfies the sharp bound
    \begin{equation}
        \label{eq:shear_bound}
        0 \le K_{ij} K^{ij} \le 3 (\kappa^*)^2;
    \end{equation}
    \item the Hamiltonian constraint \eqref{eq:wdw_hamiltonian} confines the scalar curvature to the sharp interval
    \begin{equation}
        \label{eq:curvature_bound}
        2\Lambda + 16\pi G_N\rho - 6 (\kappa^*)^2 \le {}^{(3)}R \le 2\Lambda + 16\pi G_N\rho + 2(\kappa^*)^2,
    \end{equation}
    both endpoints being attained by admissible principal curvatures.
\end{enumerate}
\end{theorem}

\begin{remark}
The theorem is conditional on $\kappa^* < \infty$. Whether a given spacetime admits Cauchy hypersurfaces with finite $\kappa^*$, whether the infimum \eqref{eq:minimax_hypersurface_problem} is attained, and whether such slices can be continued as a foliation that avoids a curvature singularity are separate questions that the theorem does not address. Near a crushing singularity $K_{ij}$ itself diverges, so a finite bound is an assumption rather than a consequence.
\end{remark}
\begin{proof}
Let $\lambda_1, \lambda_2, \lambda_3$ be the principal extrinsic curvatures of $\Sigma^*$ at point $x$, and write $a \coloneqq \kappa^*$. By definition of the $L^\infty$ operator norm, $\max_{i=1,2,3} |\lambda_i| \le a$; the minimax hypothesis imposes no further relation among the $\lambda_i$, so they range independently over $[-a,a]$. The quadratic invariant (not to be confused with the shear $\sigma_{ij}\sigma^{ij}$) satisfies:
\begin{equation}
    K_{ij} K^{ij} = \sum_{i=1}^3 \lambda_i^2 \le 3 \max_i \lambda_i^2 \le 3 a^2,
\end{equation}
proving (i); equality holds at $\lambda_1 = \lambda_2 = \lambda_3 = a$, so \eqref{eq:shear_bound} is sharp.

Substituting into the geometric Hamiltonian constraint \eqref{eq:wdw_hamiltonian} gives the exact identity:
\begin{equation}
    \begin{gathered}
    {}^{(3)}R - 2\Lambda - 16\pi G_N\rho = K_{ij} K^{ij} - K^2 \\
    = \sum_{i=1}^3 \lambda_i^2 - \Big(\sum_{i=1}^3 \lambda_i\Big)^2 = -2\, e_2(\lambda), \\
    e_2(\lambda) \coloneqq \lambda_1 \lambda_2 + \lambda_2 \lambda_3 + \lambda_3 \lambda_1.
    \end{gathered}
\end{equation}

\textbf{Tight extremization of $e_2$ on $[-a,a]^3$:} For fixed $\lambda_2, \lambda_3$, the function $e_2 = \lambda_1(\lambda_2 + \lambda_3) + \lambda_2\lambda_3$ is affine in $\lambda_1$, so on $[-a,a]$ its maximum and minimum are attained at $\lambda_1 = \pm a$; the same holds for $\lambda_2$ and $\lambda_3$. Moving one coordinate at a time to an endpoint without decreasing (resp.\ increasing) $e_2$ shows that the maximum and the minimum of $e_2$ over $[-a,a]^3$ are attained at vertices $\lambda_i = \pm a$ (they may be attained elsewhere too: the minimum $-a^2$ is attained along the whole edge $(a, -a, t)$, $t \in [-a,a]$). Using the sign symmetry $e_2(-\lambda) = e_2(\lambda)$, the vertex values reduce to exactly two cases:
\begin{align}
    \lambda_1 = \lambda_2 = \lambda_3 = a &: \quad e_2 = a^2+a^2+a^2 = 3a^2, \\
    \text{one sign flipped, e.g. } \lambda_1=\lambda_2=a,\ \lambda_3=-a &: \quad e_2 = a^2 - a^2 - a^2 = -a^2.
\end{align}
Hence $e_2(\lambda) \in [-a^2,\, 3a^2]$ for every $\lambda \in [-a,a]^3$, and both endpoints are attained.

\textbf{Conclusion:} Substituting into ${}^{(3)}R - 2\Lambda - 16\pi G_N\rho = -2e_2(\lambda)$ (which reverses the inequality direction) yields
\begin{equation}
    -6a^2 \;\le\; {}^{(3)}R - 2\Lambda - 16\pi G_N\rho \;\le\; 2a^2,
\end{equation}
establishing \eqref{eq:curvature_bound}. The lower endpoint is saturated at $\lambda_1=\lambda_2=\lambda_3=\kappa^*$ (the same configuration saturating \eqref{eq:shear_bound}), and the upper endpoint is saturated at any one-sign-flipped vertex, e.g.\ $\lambda_1=\lambda_2=\kappa^*,\ \lambda_3=-\kappa^*$.
\end{proof}

% =========================================================================
\section{Causal Jordan Embeddings and Anomaly-Free Loop Quantum Gravity}
\label{sec:jordan_lqg}
% =========================================================================

In Loop Quantum Gravity (LQG) \cite{ashtekar2004background}, physical states are expressed in terms of Wilson loops:
\begin{equation}
    \label{eq:wilson_loop_operator}
    W_\gamma(A) \coloneqq \Tr\left( \mathcal{P} \exp\left( \oint_\gamma A_a^i \tau_i \dif x^a \right) \right),
\end{equation}
where $A_a^i$ is the $SU(2)$ Ashtekar--Barbero connection.

A persistent difficulty in canonical quantization is the presence of \emph{self-intersecting loops}: when $\gamma(s_1) = \gamma(s_2)$, the action of the Hamiltonian constraint operator $\widehat{\mathcal{H}}$ at the intersection vertex is subject to the regularization ambiguities intrinsic to Thiemann's construction --- namely the dependence of the operator's action on the \emph{valence} of the vertex created, the combinatorial \emph{Mandelstam identities} relating Wilson loops through a common point (for $SU(2)$, $W_{\alpha} W_{\beta} = W_{\alpha\beta} + W_{\alpha\beta^{-1}}$), and the discretionary choice of \emph{node-splitting} prescription used to resolve the new edges.

\begin{proposition}[Timelike Curves Do Not Self-Intersect in Chronological Spacetimes]
\label{thm:jordan_chronology}
Let $(\mathcal{M}^4, g)$ be a chronological spacetime (one without closed timelike curves), for instance a stably causal one. Then every timelike curve $\gamma(\tau)$ is injective: $\gamma(\tau_1) \neq \gamma(\tau_2)$ for all $\tau_1 \neq \tau_2$.
\end{proposition}
\begin{proof}
If $\gamma(\tau_1) = \gamma(\tau_2) = p$ with $\tau_1 < \tau_2$, the restriction $\gamma|_{[\tau_1, \tau_2]}$ is a timelike curve from $p$ to $p$, so $p \in I^+(p)$, and the spacetime contains a closed timelike curve through $p$, a contradiction.
\end{proof}

\begin{remark}[Scope]
Proposition~\ref{thm:jordan_chronology} concerns \emph{timelike} curves. The loops carrying holonomies in loop quantum gravity lie in a \emph{spacelike} Cauchy surface $\Sigma$, and causality places no restriction on self-intersections of spacelike curves. Restricting to Jordan loops is therefore a choice of kinematical sector, not a consequence of chronology.
\end{remark}

\begin{conjecture}[Jordan-Loop Sector]
\label{conj:jordan_sector}
Let $\mathscr{H}_J$ be the closed subspace of the kinematical Hilbert space of loop quantum gravity spanned by spin-network states whose graph is a finite union of pairwise disjoint, smoothly embedded circles (Jordan curves) in $\Sigma$, and let $\mathscr{H}_J^\ast$ be the space of diffeomorphism-covariant distributions (elements of the algebraic dual of a dense domain) built from $\mathscr{H}_J$. There is a regularization of the Hamiltonian constraint $\widehat{\mathcal{H}}[N]$ such that
\begin{enumerate}[label=(\alph*)]
    \item $\widehat{\mathcal{H}}[N]$ does not annihilate $\mathscr{H}_J$;
    \item its (dual) action preserves $\mathscr{H}_J^\ast$;
    \item on $\mathscr{H}_J^\ast$ the commutator reproduces the hypersurface-deformation algebra,
    \begin{equation}
    [\widehat{\mathcal{H}}[N], \widehat{\mathcal{H}}[M]] = \widehat{\mathcal{H}}_a[\omega^a], \qquad \omega^a = q^{ab}(N\partial_b M - M\partial_b N),
    \end{equation}
    with the structure function quantized as an operator, and this holds for every choice of the vertex-valence and node-splitting prescriptions within the regularization.
\end{enumerate}
\end{conjecture}

\begin{remark}
Requirement (a) excludes a degenerate case. In Thiemann's regularization \cite{thiemann1998qsd} the operator $\widehat{\mathcal{H}}[N]$ acts only at vertices, through commutators with the volume operator, and the Ashtekar--Lewandowski volume vanishes at vertices of valence at most two, since it is built from triples of distinct edges. A graph made of disjoint Jordan curves has no vertices of valence greater than two, so Thiemann's $\widehat{\mathcal{H}}[N]$ annihilates $\mathscr{H}_J$. The subspace is then trivially invariant, the left-hand side of the commutation relation vanishes on it, and the right-hand side vanishes only on diffeomorphism-invariant states. A proof of the conjecture therefore needs a regularization that acts non-trivially on vertex-free graphs and does not create vertices, and closure of the quantum constraint algebra is an open problem even in restricted sectors. If Jordan curves were allowed to cross one another, the crossings would be vertices of valence four, and the sector would no longer be vertex-free; the definition of $\mathscr{H}_J$ excludes this. The heuristic motivation is that a Jordan loop has no self-intersection vertex, so the regularization choices listed above do not arise on the initial state.
\end{remark}

% =========================================================================
\section{Tensor Network Holography and the Emergent Einstein Equations}
\label{sec:holographic_emergence}
% =========================================================================

We now link the discrete simplicial theory and minimax foliations to the continuous tensor network framework established in \cite{silvafilho2026functional, silvafilho2026flows, silvafilho2026spacetime, verstraete2010continuous, swingle2012entanglement}.

\begin{theorem}[Wald Symplectic Equivalence: Linearized Einstein Equations and Second-Order Extension (after \cite{faulkner2014gravitation, faulkner2017nonlinear})]
\label{thm:grand_einstein_emergence}
Let $|\Psi\rangle$ be a state in a continuous tensor network (cMERA/cMPS) which is \emph{assumed} to admit a semiclassical bulk dual $g_{\mu\nu} = g_{\mu\nu}^{\mathrm{AdS}} + h_{\mu\nu}$, perturbatively close to pure AdS, in which the Ryu--Takayanagi formula holds for ball-shaped boundary regions, with $h_{\mu\nu}$ treated to the order indicated below. (For holographic CFTs these hypotheses are those of \cite{faulkner2014gravitation, faulkner2017nonlinear}; that a given cMERA/cMPS satisfies them is an assumption, not a result of this treatise.)
\begin{enumerate}[label=(\roman*)]
    \item \textbf{First order.} The first variation of the boundary first law of entanglement, $\delta S_A = \delta \langle H_A \rangle$, holding for all boundary spherical regions $A$, is equivalent via the Wald symplectic identity $\int_{\partial \Sigma} \chi = \int_\Sigma \dif \chi$ to the \emph{linearized} vacuum Einstein equation
    \begin{equation}
        \delta\!\left( G_{\mu\nu} + \Lambda g_{\mu\nu} \right) = 0, \quad \forall (x, z) \in M,
    \end{equation}
    i.e.\ $h_{\mu\nu}$ satisfies the linearized Einstein equation to first order in the perturbation \cite{faulkner2014gravitation}. With bulk matter, Swingle and Van Raamsdonk \cite{swingle2014universality} argue, for semiclassical bulk states and with the quantum corrections to the entanglement formula, that the source term $-8\pi G_N \delta\langle T_{\mu\nu}^{\mathrm{matter}}\rangle$ appears in the bracket.
    \item \textbf{Second order.} The second-order variation of the relative entropy $S(\rho_A \| \sigma_A)$ of ball-shaped regions equals the canonical energy of the bulk perturbation \cite{lashkari2016canonical}; combined with the holographic entanglement relation at second order, this yields the Einstein equations at second order, as shown in \cite{faulkner2017nonlinear} for the class of states prepared by sources in the Euclidean path integral, to second order in the sources. Positivity of relative entropy then implies positivity of canonical energy, i.e.\ a stability condition, but does not by itself fix the higher-order terms.
\end{enumerate}
\end{theorem}
\begin{proof}
Both items are imported from the cited works; see also Section~3 of \cite{silvafilho2026spacetime}. At first order, the Noether charge $(d-1)$-form on the extremal surface $\gamma_A$ integrates to $\delta \Area(\gamma_A)/(4G_N) = \delta S_A$, while on the boundary it integrates to the modular Hamiltonian $\delta\langle H_A\rangle$; the identity $\dif\chi = -2\xi^\mu \delta(G_{\mu\nu} + \Lambda g_{\mu\nu})\mathbf{\epsilon}^\nu$ forces the linearized combination to vanish pointwise for all balls $A$, by the sweeping argument of Theorem~3.2 of \cite{silvafilho2026spacetime}. This argument is intrinsically linear: it uses only the first variation of $S_A$, and therefore fixes $h_{\mu\nu}$ only to $\mathcal{O}(h)$. Extending the identification to $\mathcal{O}(h^2)$ uses the second-order analysis of \cite{faulkner2017nonlinear} (item (ii)), under the same holographic hypotheses. Full non-perturbative (all-orders) equivalence to the non-linear Einstein equations is not established by this argument.
\end{proof}

\begin{conjecture}[Mean Curvature Flow as a Search Procedure for Ryu--Takayanagi Surfaces]
\label{thm:rt_mcf_convergence}
Call a hypersurface of the emergent bulk geometry an \emph{admissible cut} if it is anchored on $\partial A$ and homologous to $A$, and measure areas with the cutoff at distance $\epsilon$ from the asymptotic boundary, subtracting the $\epsilon$-divergent part (which depends only on $\partial A$). The mean curvature flow $\partial_t \gamma = \mathbf{H}_\gamma$ with boundary fixed at $\partial A$, in the level-set sense, started from any admissible cut, exists for all time and converges to an admissible extremal cut ($\mathbf{H} = 0$). The Ryu--Takayanagi surface $\gamma_A$ \cite{ryu2006holographic} is the extremal cut of least regularized area among these limits.
\end{conjecture}

\begin{remark}
Along any smooth portion of the flow the area is non-increasing, since (first variation of area)
\[
\frac{\dif}{\dif t}\Area(\gamma(t)) = -\int_{\gamma(t)} \|\mathbf{H}\|^2 \dif\Area \le 0 .
\]
The conjecture does not assert that every initial cut flows to $\gamma_A$, and it could not: mean curvature flow is the gradient flow of area, so every extremal cut is a fixed point, minimizing or not. In the time slice $\mathbb{H}^2$ of AdS$_3$ (radius $L$), let $A = [0,1] \cup [1+\delta, 2+\delta]$ on the boundary. The pair of geodesics over the two intervals and the pair over $[0, 2+\delta]$ and $[1, 1+\delta]$ are both admissible and have $\mathbf{H} = 0$. Their regularized lengths are $4L\log(1/\epsilon)$ and $2L\log((2+\delta)/\epsilon) + 2L\log(\delta/\epsilon)$; the second is smaller exactly when $\delta(2+\delta) < 1$. For $\delta = 0.1$ the first pair exceeds the second by $-2L\log 0.21 \approx 3.12\,L$, yet the flow started from it does not move. The flow can also form singularities, which is why the level-set formulation is used. Identifying the area of $\gamma_A$ with $4 G_N S_A$ needs the Ryu--Takayanagi formula, which is an input here and not a consequence. The same statement appears as Conjecture~4.5 of Chapter~11 \cite{silvafilho2026spacetime}.
\end{remark}

\begin{remark}[Fefferman--Graham coefficients and relative entropy]
\label{rem:fefferman_graham}
In Fefferman--Graham gauge an asymptotically AdS$_{d+1}$ metric reads
\begin{equation}
    \label{eq:fefferman_graham_exact}
    g = \frac{L^2}{z^2}\left[ \dif z^2 + \left( g_{(0)\mu\nu} + z^2 g_{(2)\mu\nu} + \dots + z^d g_{(d)\mu\nu} + \dots \right) \dif x^\mu \dif x^\nu \right].
\end{equation}
The coefficients $g_{(2k)}$ with $2k < d$ are local functionals of the boundary metric $g_{(0)}$; for instance, for $d > 2$,
\[
g_{(2)\mu\nu} = -\frac{1}{d-2}\left( R_{\mu\nu}[g_{(0)}] - \frac{R[g_{(0)}]}{2(d-1)}\, g_{(0)\mu\nu} \right).
\]
The state enters first at order $z^d$: $g_{(d)\mu\nu} = \frac{16\pi G_N}{d\,L^{d-1}}\langle T_{\mu\nu}\rangle + X_{\mu\nu}[g_{(0)}]$, where $X_{\mu\nu}$ is a local term that vanishes for flat $g_{(0)}$ and for odd $d$ \cite{deharo2001holographic}. Relative entropy enters through Theorem~\ref{thm:grand_einstein_emergence}: its first variation, the first law, constrains $\delta g_{(d)}$, and its second derivative $\mathcal{C}_2 = \frac{\dif^2}{\dif\lambda^2} S(\rho_A(\lambda) \| \rho_A(0))|_{\lambda=0}$, the relative-entropy (Kubo--Mori) Fisher information of the perturbation, equals the canonical energy \cite{lashkari2016canonical}. It is not a Fefferman--Graham coefficient. An identification $g_{(2)\mu\nu} = \mathcal{C}_2/L^2$ would fail: for a flat boundary $g_{(2)} = 0$ for every state, whereas $\mathcal{C}_2 > 0$ for every non-trivial perturbation (numerical check in the accompanying code repository). Whether the higher coefficients of the expansion of relative entropy in the perturbation constrain the nonlinear terms of $g_{(d)}$ and beyond, in the sense of an all-orders version of Theorem~\ref{thm:grand_einstein_emergence}, is open.
\end{remark}

% =========================================================================
\section{The Graphon Ricci Condensation and Topology Transitions}
\label{sec:graphon_condensation}
% =========================================================================

Can a discrete, foam-like connectivity structure condense into a smooth geometry? This section records what the graphon Ricci flow of Chapter~2 does and does not show about that question.

In \cite{silvafilho2026flows}, we formulated the continuous Graphon Ricci Flow on graph kernels $W \in L^\infty([0, 1]^2)$:
\begin{equation}
    \label{eq:graphon_ricci_unified}
    \partial_t W(t, x, y) = 2 \kappa_W(t, x, y) W(t, x, y),
\end{equation}
where $W$ is a connection strength (not a distance), so entries with negative curvature decrease. With the opposite sign, negatively curved bottlenecks would grow.

\begin{proposition}[Exponential Decay of Negatively Curved Entries]
\label{thm:neckpinch_condensation}
Let $W(t)$ be a solution of \eqref{eq:graphon_ricci_unified} on $[0,T]$ with $W_0 = W(0) \ge 0$, and let $B \subset [0,1]^2$ be a set of positive measure on which $\kappa_{W(s)}(x,y) \le -c < 0$ for all $s \in [0,T]$. Then for $(x,y) \in B$ and $t \in [0,T]$,
\begin{equation}
    0 \le W(t,x,y) \le W_0(x,y)\, e^{-2ct},
\end{equation}
so for every $\delta \in (0, \|W_0\|_{L^\infty(B)})$ the entries on $B$ are at most $\delta$ at all times $t \in [t_\delta, T]$, where $t_\delta = \frac{1}{2c}\ln\frac{\|W_0\|_{L^\infty(B)}}{\delta}$ (provided $t_\delta \le T$).
\end{proposition}
\begin{proof}
For fixed $(x,y)$, \eqref{eq:graphon_ricci_unified} is the linear ODE $\dot W = 2\kappa_W W$, whose solution is $W(t) = W_0 \exp(2\int_0^t \kappa_W(s) \dif s)$. The bound $\kappa_W \le -c$ gives $2\int_0^t \kappa_W \le -2ct$, hence the estimate.
\end{proof}

\begin{remark}[Scope of the curvature hypothesis]
Ollivier's coarse Ricci curvature $\kappa_W(x,y) = 1 - W_1(m_x, m_y)/d(x,y)$ depends on the metric $d$ on $[0,1]$ used in the transport distance $W_1$, and Chapter~2 leaves $d$ unspecified, since a graphon carries no canonical graph distance. The admissible range of $c$ depends on this choice. With the discrete metric ($d(x,y) = 1$ for $x \ne y$), $W_1$ is the total-variation distance and $\kappa_W \ge 0$, so the hypothesis $\kappa_W \le -c < 0$ never holds and the proposition is vacuous. With the Euclidean metric $|x - y|$ the curvature is not bounded below: for the symmetric block kernel equal to $1$ on $I_1 \times I_2 \cup I_3 \times I_4$ and its transpose, with $I_j$ the quarters of $[0,1]$ and a small constant background elsewhere, two nearby points $x < y$ on either side of a block boundary have measures $m_x, m_y$ at distance of order one while $d(x,y) \to 0$, and one finds $\kappa_W(0.49, 0.51) \approx -36$ and $\kappa_W(0.499, 0.501) \approx -372$ (numerical check in the accompanying code repository). The lower bound $\kappa \ge -2$ holds for the graph distance of a graph with unit-length edges, where neighbouring measures are supported at distance at most one from their centres; it has no analogue for general $d$. Whether the hypothesis $\kappa_W \le -c$ on $B$ holds, and is preserved along the flow, therefore has to be checked for each choice of $d$ and each initial kernel.
\end{remark}

\begin{conjecture}[Graphon Condensation]
Starting from a random kernel $W_0$ modelling a branched-polymer phase, the graphon Ricci flow disconnects one-dimensional bottlenecks in finite time, and the remaining positively curved clusters converge, after suitable rescaling, to the kernel of a smooth $4$-dimensional geometry.
\end{conjecture}

\begin{remark}
Convergence in the cut metric $\delta_\square$ is convergence of graphons, i.e.\ of symmetric measurable kernels. A graphon limit carries no dimension, differentiable structure or metric tensor, so the phrase ``converges to a $4$-dimensional Einstein manifold'' has no meaning until a reconstruction map from kernels to Riemannian manifolds is specified. Constructing such a map is part of the conjecture.
\end{remark}

% =========================================================================
\section{Quantum Chaos, Black Hole Horizons, and Kac--Rice Morse Stratification}
\label{sec:chaos_kac_rice}
% =========================================================================

At black hole event horizons, quantum information scrambling reaches the universal thermal Lyapunov bound of Maldacena, Shenker, and Stanford (MSS) \cite{maldacena2016bound}:
\begin{equation}
    \label{eq:mss_lyapunov_bound}
    \lambda_L \le \frac{2\pi k_B T}{\hbar}.
\end{equation}

Two facts frame this section. First, the Sachdev--Ye--Kitaev (SYK) model of $N$ Majorana fermions with random $q$-body couplings has an extensive zero-temperature entropy $S_0 N$ and an exponent $\lambda_L$ that approaches the bound \eqref{eq:mss_lyapunov_bound} only in the strong-coupling limit $\beta J \to \infty$; at finite coupling $\lambda_L < 2\pi/\beta$ \cite{maldacena2016remarks}. Second, the Kac--Rice count of critical points of a random spherical $k$-spin landscape is a different quantity. For the isotropic model with i.i.d.\ Gaussian couplings (Theorem~5.6 of \cite{silvafilho2026functional}, where the dimension and the degree are written $(n, d)$ instead of $(N, k)$), $\frac1N \log \mathbb{E}[\mathcal{N}_{\mathrm{crit}}] \to \theta(k) = \frac12\log(k-1)$ for $k \ge 3$, a rate due to Auffinger, Ben~Arous and \v{C}ern\'y \cite{auffinger2013random}, while for $k = 2$ there are exactly $2N$ critical points almost surely. This complexity counts critical points of a classical energy function on $\Sph^{N-1}$, not quantum microstates.

\begin{conjecture}[Kac--Rice Complexity and Horizon Entropy]
\label{thm:horizon_morse_saturation}
There is a class of random tensor Hamiltonians for which the annealed complexity $\theta(k)$ of the classical landscape coincides, to leading order in $N$, with the thermodynamic entropy of a dual near-extremal black hole, and in which the large-$N$ out-of-time-order correlator $F(t) = 1 - \frac{C}{N} e^{\lambda_L t} + \mathcal{O}(N^{-2})$ has $\lambda_L \to 2\pi/\beta$ in an appropriate strong-coupling limit.
\end{conjecture}

\begin{remark}
Neither half of the conjecture follows from Kac--Rice counting or from sub-Gaussian concentration. Concentration controls fluctuations of the landscape; it does not determine the Lyapunov exponent, and it does not identify critical points with microstates.
\end{remark}


% =========================================================================
\section{Phenomenological Geometric Parametrization of Fundamental Interactions}
\label{sec:geometric_condensation_interactions}
% =========================================================================

This section is of a different kind from the preceding ones: it contains no theorems. A unified framework often seeks to relate the standard model of particle physics to low-energy geometric condensations. In this section, we present a \emph{phenomenological model} and \emph{geometric parametrization}—rather than a parameter-free first-principles derivation—illustrating how a schematic action $\mathcal{S}_{\mathrm{univ}}$ (specified below only through its gauge-kinetic form and its Dirac operator; we do not define it completely) and the Dirac operator over the simplicial product $\Delta_4 \times \Delta_2$ can be structurally aligned with known particle hierarchies. We explicitly emphasize that this geometric mapping uses empirical Standard Model parameters (such as the strong coupling $\alpha_s(M_Z)$ and the electroweak vacuum expectation value $v$) as inputs to calibrate the geometric embedding scales, rather than claiming an unconstrained, parameter-free prediction.

\subsection{The Universal Action and Phenomenological Yukawa Coupling}
The action $\mathcal{S}_{\mathrm{univ}}$ must couple the gauge fields, the fermions $\boldsymbol{\Psi}$, and the scalar Higgs field $\Phi \in \Gamma(\mathcal{H})$. To ensure strict coercivity and a positive-definite norm, the bilinear form over the gauge algebra is evaluated in the Euclidean compact rotation algebra $\mathfrak{so}(4) \cong \mathfrak{su}(2) \oplus \mathfrak{su}(2)$ rather than the non-compact Lorentz algebra $\mathfrak{so}(3,1)$:
\begin{equation}
    \label{eq:bilinear_form_su2}
    \langle \boldsymbol{\Omega}, \boldsymbol{\Omega} \rangle_{\mathcal{G}} = \frac{1}{g_3^2}\Tr_c(G \wedge \star_{\mathcal{G}} G) + \frac{1}{g_2^2}\Tr_L(W \wedge \star_{\mathcal{G}} W) + \frac{1}{g_1^2}(B \wedge \star_{\mathcal{G}} B) + \frac{1}{g_{\mathrm{grav}}^2}\Tr_E(\mathcal{R} \wedge \star_{\mathcal{G}} \mathcal{R}).
\end{equation}
For the product $\mathcal{M} = \Delta_4 \times \Delta_2$ we take the Dirac operator in the standard form of a product of spectral triples \cite{chamseddine2007gravity}:
\begin{equation}
    \label{eq:dirac_weingarten}
    \boldsymbol{\mathcal{D}}_{\mathcal{M}}(\Phi) = \boldsymbol{\mathcal{D}}_{\Delta_4} \otimes \mathbf{1}_{\Delta_2} + \gamma_5 \otimes \boldsymbol{\mathcal{W}}_{\Delta_2}(\Phi),
\end{equation}
where $\gamma_5$ is the chirality grading of the four-dimensional factor, which anticommutes with $\boldsymbol{\mathcal{D}}_{\Delta_4}$, and $\boldsymbol{\mathcal{W}}_{\Delta_2}(\Phi) = \mathbf{Y}_{\mathrm{circ}} \Phi(x)$ is the Weingarten endomorphism coupled to the Higgs field. Because of the anticommutation the cross terms cancel, $\boldsymbol{\mathcal{D}}_{\mathcal{M}}^2 = \boldsymbol{\mathcal{D}}_{\Delta_4}^2 \otimes \mathbf{1} + \mathbf{1} \otimes \boldsymbol{\mathcal{W}}^2$, and a mode of momentum $k$ in a generation with eigenvalue $m$ of $\boldsymbol{\mathcal{W}}$ has $\boldsymbol{\mathcal{D}}_{\mathcal{M}}$-eigenvalues $\pm\sqrt{k^2 + m^2}$, i.e.\ a mass gap. With $\mathbf{1}_{\Delta_4}$ in place of $\gamma_5$ the two terms commute and the eigenvalues would be $\pm|k| + m$, a shift of the spectrum without a gap. Upon spontaneous symmetry breaking $\langle \Phi \rangle = v / \sqrt{2}$, this gives the fermion mass matrix $\mathbf{M}_f = \frac{v}{\sqrt{2}} \mathbf{Y}_{\mathrm{circ}}$. The factor $\Delta_4$ is a simplex with boundary, and a self-adjoint Dirac operator on it requires boundary conditions, which we do not specify; the statements above concern only the algebraic structure of \eqref{eq:dirac_weingarten}.

The circulant form of $\mathbf{Y}_{\mathrm{circ}}$ is an assumption. Gauge invariance allows any complex $3 \times 3$ Yukawa matrix. We make it precise as follows. Let $P$ be the cyclic permutation matrix of the three generations and
\begin{equation}
    \label{eq:circulant_root}
    \mathbf{C} = a\,\mathbf{1} + b\,e^{i\delta} P + b\,e^{-i\delta} P^{-1}, \qquad a, b \ge 0,
\end{equation}
a Hermitian circulant matrix with eigenvalues $v_j = a + 2b\cos(\delta + 2\pi j/3)$, $j = 1,2,3$. We set $\mathbf{Y}_{\mathrm{circ}} = (\sqrt2/v)\,\mathbf{C}^2$. Then $\mathbf{M}_f = \mathbf{C}^2$ is again circulant, and its eigenvalues are the masses $m_j = v_j^2$. The parametrization $a + 2b\cos(\cdot)$ therefore describes the \emph{square roots} of the masses, as in the Koide relation below, and not the masses themselves. Every circulant matrix is diagonalized by the discrete Fourier matrix $F_{jk} = \omega^{jk}/\sqrt3$, $\omega = e^{2\pi i/3}$, whatever $(a, b, \delta)$ are. So if the up-type and the down-type mass matrices are both of this form in the same basis, one unitary diagonalizes both, and the CKM matrix is a permutation matrix up to phases. The circulant ansatz thus gives no quark mixing, and in particular no Cabibbo angle. Mixing would need an additional, non-circulant structure in at least one quark sector, which we do not supply.
\subsection{Algebraic Symmetries of the Fermion Generation Simplex}
The geometric reduction of the fermion generation parameter space is governed by the permutation symmetry $S_3$ acting on the generation simplex $\Delta_2$. For the lepton sector, the representation theory of $S_3$ motivates the bidirectional equivalence of quadratic norm equipartition, recovering the empirical Koide mass relation \cite{koide1981}:
\begin{equation}
    \|\mathbf{v}_{\mathbf{2}}\|^2 = \|\mathbf{v}_{\mathbf{1}}\|^2 \iff 6b^2 = 3a^2 \iff \frac{b}{a} = \frac{1}{\sqrt{2}} \iff Q_l \equiv \frac{\Tr(\mathbf{M}_l)}{\Tr(\sqrt{\mathbf{M}_l})^2} = \frac{2}{3},
\end{equation}
where $\mathbf{v} = (\sqrt{m_e}, \sqrt{m_\mu}, \sqrt{m_\tau})$ is decomposed as $\mathbf{v} = \mathbf{v}_{\mathbf{1}} + \mathbf{v}_{\mathbf{2}}$ into its component along $(1,1,1)$ (the trivial representation of $S_3$) and its component in the orthogonal plane (the two-dimensional standard representation), written as $v_j = a + 2b\cos(\delta + 2\pi j/3)$, $j = 1,2,3$, with $a, b \ge 0$ and a phase $\delta$; then $\|\mathbf{v}_{\mathbf{1}}\|^2 = 3a^2$, $\|\mathbf{v}_{\mathbf{2}}\|^2 = 6b^2$ and $Q_l = 1/3 + 2b^2/(3a^2)$.

The equivalence above is algebraic: $Q_l = 2/3$ is the statement that the vector $(\sqrt{m_e},\allowbreak\sqrt{m_\mu},\allowbreak\sqrt{m_\tau})$ makes an angle of $45^\circ$ with $(1,1,1)$, a known reformulation of Koide's empirical relation. The $S_3$ structure makes the relation natural but does not derive it.

For the heavy quarks $(c, b, t)$ no comparable statement is available. The value of $Q_{cbt}$ is scheme-dependent: with $\overline{\mathrm{MS}}$ masses at $M_Z$ one finds $Q_{cbt} \approx 0.72$, with pole masses $Q_{cbt} \approx 0.65$. The expression $\frac{2}{3}\left( 1 + \alpha_s(M_Z)/\sqrt{3} \right) \approx 0.712$ lies in this range, but it is a numerical fit with $\alpha_s(M_Z)$ as input and no derivation: the factor $1/\sqrt{3}$ is neither the colour Casimir $C_F = 4/3$ nor a known one-loop coefficient. Given the spread between schemes, agreement with it carries no evidential weight.

\subsection{Higgs Mass and Parameter Count}
Starting from the reference value $m_H^{(0)} = v/2 = 123.11\text{ GeV}$ ($\lambda = 1/8$ with $m_H = v\sqrt{2\lambda}$), a shift of the quartic coupling $\Delta\lambda \approx +0.0043$ yields the measured value $m_H \approx 125.2\text{ GeV}$. The shift $\Delta\lambda$ is chosen to reproduce the measured mass, so this is a one-parameter fit to one observable and carries no predictive content. In total, the parametrization of this section uses at least as many inputs ($v$, $\alpha_s(M_Z)$, $\Delta\lambda$, and the lepton masses entering the Koide relation) as the observables it is compared with. We therefore present it as a consistency check, not a prediction.

We do not currently have a geometric expression for $J_{\mathrm{CP}}$, whose measured value is $J_{\mathrm{CP}} \approx 3.1 \times 10^{-5}$.

% =========================================================================
\section{Analytical Solutions in Truncated Models}
\label{sec:analytical_solutions}
% =========================================================================

We collect closed-form solutions of specific reduced or truncated toy models suggested by the preceding sections. Each is exact within its model. None is a solution of the full theory.

\subsection{A One-Dimensional Closed-Loop Toy Model}
Motivated by Conjecture~\ref{conj:jordan_sector}, we consider the \emph{minisuperspace truncation} of the full Wheeler--DeWitt theory obtained by restricting attention to a single closed-loop mode $\gamma(s)$ and freezing all other kinematical degrees of freedom; this is the standard reduction used to extract a tractable one-dimensional model from the infinite-dimensional constraint, and it should not be read as a solution of the full multi-loop Hamiltonian constraint. We \emph{postulate} that, under the extrinsic curvature bound $\|\II\|_{L^\infty} = \kappa^*$ (Theorem~\ref{thm:minimax_hamiltonian_regularization}), the constraint restricted to this single-mode sector takes the one-dimensional Sturm--Liouville form (this reduction is not derived here):
\begin{equation}
    \label{eq:wdw_exact_sl}
    \left[ -\frac{\hbar^2}{\kappa^*} \frac{\dif^2}{\dif s^2} + V_{\mathrm{ADM}}(s) \right] \psi_\gamma(s) = 0.
\end{equation}
When the spatial curvature potential is linearized near the boundary of the embedding $V_{\mathrm{ADM}}(s) = V_0 + \kappa^* \sigma s$, equation \eqref{eq:wdw_exact_sl} has the exact analytical solution:
\begin{equation}
    \label{eq:wdw_airy_solution}
    \psi_\gamma(s) = c_1 \operatorname{Ai}\left( \left(\frac{(\kappa^*)^2 \sigma}{\hbar^2}\right)^{1/3} \left( s + \frac{V_0}{\kappa^* \sigma} \right) \right) + c_2 \operatorname{Bi}\left( \left(\frac{(\kappa^*)^2 \sigma}{\hbar^2}\right)^{1/3} \left( s + \frac{V_0}{\kappa^* \sigma} \right) \right),
\end{equation}
where $\operatorname{Ai}$ and $\operatorname{Bi}$ are the standard Airy functions. On a half-line $s \in [s_0,\infty)$ with $\sigma > 0$, requiring decay as $s \to \infty$ enforces $c_2 = 0$, and a Dirichlet condition $\psi_\gamma(s_0) = 0$ then restricts, for fixed $s_0$, the constant $V_0$ to a discrete set determined by the zeros of $\operatorname{Ai}$ (the equation has no separate eigenvalue parameter). The units of $\hbar^2/\kappa^*$ and $\kappa^*\sigma$ are fixed by the postulated form of \eqref{eq:wdw_exact_sl}, not derived. On a closed loop the parameter $s$ is periodic, and the linearized potential is not. Equation \eqref{eq:wdw_airy_solution} is therefore only a local solution, valid in the region where the linearization holds; it does not by itself determine a spectrum on the loop.

\subsection{Closed-Form Dynamical Formula for the Running Spectral Dimension}
In Causal Dynamical Triangulations in $3+1$ dimensions \cite{ambjorn2005spectral}, the reduction $d_s = 4 \to 2$ was found numerically. In $2+1$ dimensions, Sotiriou, Visser and Weinfurtner \cite{sotiriou2011spectral} showed that the flow of the spectral dimension in CDT is reproduced to high accuracy by diffusion with Ho\v{r}ava--Lifshitz-type anisotropic dispersion relations. Here we consider a different, isotropic operator: for the \emph{assumed} two-scale symbol $k^2(1 + \ell_P^2 k^2)$ in four Euclidean dimensions (Proposition~\ref{thm:running_spectral_dimension} with $\alpha = 1$ and $\alpha = 2$ combined), the return probability is an elementary integral that can be evaluated in closed form. (For the anisotropic operator $\omega^2 + k^2 + \ell_P^2 k^4$ in $3+1$ dimensions the ultraviolet value would be $d_s = 1 + 3/2 = 5/2$, not $2$.) Writing $u = k^2$ and $z \coloneqq \sqrt{\tau}/(2\ell_P)$, the radial momentum integral retains the full polynomial prefactor $u$:
\begin{equation}
    P(\tau; \mathbf{x}, \mathbf{x}) = \frac{1}{(2\pi)^4} \int_{\R^4} e^{-\tau (k^2 + \ell_P^2 k^4)} \dif^4 k = \frac{1}{32 \pi^2 \ell_P^2 \tau} \left[ 1 - \sqrt{\pi}\, z\, e^{z^2} \operatorname{erfc}(z) \right].
\end{equation}
Taking the logarithmic derivative $d_s(\tau) \coloneqq -2 \frac{\dif \ln P(\tau)}{\dif \ln \tau}$ yields the closed-form formula:
\begin{equation}
    \label{eq:analytical_ds_formula}
    d_s(\tau) = 1 - \frac{\tau}{2\ell_P^2} + \left[ 1 - \frac{\sqrt{\pi \tau}}{2\ell_P} \exp\!\left(\frac{\tau}{4\ell_P^2}\right) \operatorname{erfc}\!\left(\frac{\sqrt{\tau}}{2\ell_P}\right) \right]^{-1}.
\end{equation}
Asymptotic expansion confirms:
\begin{equation}
    \lim_{\tau \to \infty} d_s(\tau) = 4 \quad (\text{Macroscopic Classical IR}), \qquad \lim_{\tau \to 0} d_s(\tau) = 2 \quad (\text{Planckian Quantum UV}),
\end{equation}
where in the infrared limit the $-\tau/(2\ell_P^2) = -2z^2$ term cancels against the $2z^2 + 3 + \mathcal{O}(z^{-2})$ reciprocal tail, restoring $d_s \to 4$; the approach is slow, $d_s = 4 - 12\ell_P^2/\tau + \mathcal{O}(\ell_P^4/\tau^2)$. Closed forms of this kind are known: Sotiriou, Visser and Weinfurtner \cite{sotiriou2011dispersion} evaluate the return probability of the anisotropic operator $-\partial_t^2 + f(-\nabla^2)$ with the two-scale symbol $f(k^2) = k^2(1 + k^2/4K^2)$ in terms of the error function in $2+1$ dimensions, and with $f(k^2) = k^2(1 + k^2/8K^2)$ in terms of Bessel functions in $3+1$ dimensions. The four-dimensional isotropic case reduces to theirs: with $I_n(\tau) = \int_0^\infty u^n e^{-\tau(u + \ell_P^2 u^2)} \dif u$, integration by parts gives $I_1 = (\tau^{-1} - I_0)/(2\ell_P^2)$, and $I_0$ is, up to a constant factor, their two-dimensional spatial integral. We have checked \eqref{eq:analytical_ds_formula} against direct numerical quadrature of the return probability for $\tau/\ell_P^2 \in [10^{-3}, 10^{2}]$. The formula gives the spectral dimension of this particular diffusion model. Whether quantum spacetime diffuses according to this symbol is an assumption.

\subsection{The Ryu--Takayanagi Hemisphere as a Stationary Point of Mean Curvature Flow}
In the hyperbolic metric $g = \frac{L^2}{z^2}(\dif z^2 + \dif \mathbf{x}^2)$, $\mathbf{x} \in \R^{d-1}$, $z>0$, consider the ball $A = \{ \|\mathbf{x}\| \le R_0 \}$ on the boundary $z = 0$ and the hemisphere
\begin{equation}
    \gamma_A = \{ z^2 + \|\mathbf{x}\|^2 = R_0^2 \}.
\end{equation}
Spheres centred on the boundary meet it orthogonally and are totally geodesic in the half-space model. Hence $\II_{\gamma_A} = 0$, in particular $\mathbf{H}_{\gamma_A} = 0$, and $\gamma_A$ is the Ryu--Takayanagi surface of $A$ \cite{ryu2006holographic}. It is a \emph{stationary} solution of the mean curvature flow, so along the flow started at $\gamma_A$ the entropy $S_A$ is constant. In particular there is no shrinking self-similar solution $R(t) = R_0 e^{-t/L}$ through hemispheres. The area of $\gamma_A$ diverges near $z = 0$; with a cutoff $z \ge \epsilon$ it reproduces the standard divergence of $S_A$, proportional to $\epsilon^{-(d-2)}$ (an area law) for $d \ge 3$ and logarithmic in $\epsilon$ for $d = 2$.

\subsection{Disconnection Time in the Graphon Ricci Flow}
Under the hypothesis of Proposition~\ref{thm:neckpinch_condensation} ($\kappa_W \le -c < 0$ on the bottleneck set $B$), and with $W_{\mathrm{split}}$ the kernel obtained by setting $W$ to zero on $B$, one has $\cutnorm{W(t) - W_{\mathrm{split}}} \le \|W(t)\mathbf{1}_B\|_{L^1} \le \|W_0\|_{L^\infty} |B|\, e^{-2ct}$. This falls below a threshold $\delta$ by the time
\begin{equation}
    \label{eq:exact_surgery_time}
    T_{\mathrm{surgery}} = \frac{1}{2 c} \ln\left( \frac{\|W_0\|_{L^\infty}|B|}{\delta} \right).
\end{equation}
This controls only the disconnection of $B$. Whether the complement evolves toward a smooth geometry is the content of the Graphon Condensation conjecture above.

\subsection{Thin-Shell Profiles of Constant Extrinsic Curvature}
Consider a timelike thin shell $r = r(\ell)$ in the Israel formalism, with $\ell$ the proper time along the shell and a Schwarzschild region of mass $M$ on one side, and \emph{impose} that the angular extrinsic curvature on that side be constant, $K^\theta{}_\theta = \sqrt{1 - 2M/r + (\dif r/\dif\ell)^2}\,/\,r \equiv \kappa^*$. (Chebyshev equioscillation of Chapter~7 motivates this choice. It does not imply it for this problem.) With this ansatz the equation of motion of the shell takes the form
\begin{equation}
    \left( \frac{\dif r}{\dif \ell} \right)^2 + 1 - \frac{2 M}{r} = (\kappa^*)^2 r^2,
\end{equation}
which is integrated by the quadrature
\begin{equation}
    \label{eq:jacobi_throat_solution}
    \ell(r) = \int_{r_{\min}}^{r} \frac{\dif r'}{\sqrt{(\kappa^*)^2 r'^2 - 1 + 2M/r'}},
\end{equation}
an elliptic integral, since the radicand times $r'$ is the cubic $(\kappa^*)^2 r'^3 - r' + 2M$. This cubic has its minimum over $r' > 0$ at $r' = 1/(\sqrt3\,\kappa^*)$, where it equals $2M - 2/(3\sqrt3\,\kappa^*)$. Hence a turning point $r_{\min} > 0$ (the largest positive root) exists if and only if $27(\kappa^*)^2M^2 < 1$; otherwise the radicand is positive for all $r' > 0$ and the shell has no bounce. Whether a given profile requires exotic matter, and how much, follows from the Israel stress tensor $S_{ab}$. That requirement is not minimized by this construction in any proved sense.

% =========================================================================
\section{Observational Signatures and Experimental Testability}
\label{sec:experimental}
% =========================================================================

A framework that aims at quantum gravity should state what it implies for observation, including the case where it implies nothing observable. This section gives order-of-magnitude estimates; the details, and the script that computes them, are in Chapter~13 and in the accompanying code repository. We use the non-reduced Planck mass $M_P = \sqrt{\hbar c/G} \approx 1.22 \times 10^{19}$\,GeV; with the reduced mass $M_P/\sqrt{8\pi}$ the numbers change by factors of order $8\pi$.

\subsection{Gravitational-Wave Dispersion}
Suppose, as a phenomenological ansatz not derived from the preceding sections, that the running spectral dimension is accompanied by a modified graviton dispersion relation with a scale $\ell_\ast$:
\begin{equation}
    \label{eq:graviton_dispersion}
    \omega^2(k) = c^2 k^2 \left( 1 + \xi \ell_\ast^2 k^2 \right).
\end{equation}
The diffusion model of Section~\ref{sec:analytical_solutions} does not fix $\xi$. Continued to Lorentzian signature, the isotropic symbol $p^2(1 + \ell_P^2 p^2)$ vanishes on $p^2 = 0$, a massless branch without dispersion ($\xi = 0$), and on $p^2 = -1/\ell_P^2$, a massive mode with negative residue (a ghost). The anisotropic operator $\omega^2 + c^2(k^2 + \ell_P^2 k^4)$, which would give $\xi = 1$, has ultraviolet spectral dimension $5/2$ rather than $2$ (Section~\ref{sec:analytical_solutions}). We therefore treat $\xi$ as a free parameter of order unity; the estimates are linear in $\xi$, and we quote them for $\xi = 1/2$, the value used in Chapter~13. The group velocity is $v_g \approx c(1 + \frac{3}{2}\xi \ell_\ast^2 k^2)$. A graviton observed at frequency $f$ had physical wavenumber $2\pi f(1+z')/c$ at redshift $z'$, so for a source at redshift $z$ two frequencies $f_1 < f_2$ emitted together arrive separated by \cite{jacob2008lorentz, mirshekari2012constraining}
\begin{equation}
    |\Delta t_{\mathrm{disp}}| = \frac{6 \pi^2 \xi \ell_\ast^2}{c^3}\, D_2(z) \left( f_2^2 - f_1^2 \right), \qquad D_2(z) = \frac{c}{H_0} \int_0^z \frac{(1+z')^2 \dif z'}{\sqrt{\Omega_m(1+z')^3 + \Omega_\Lambda}}.
\end{equation}
For $\ell_\ast = \ell_P$, $\xi = 1/2$, $f \in [10, 10^3]$\,Hz, $H_0 = 70$\,km\,s$^{-1}$\,Mpc$^{-1}$, $\Omega_m = 0.3$ and redshifts $z = 1$--$8$ this gives $|\Delta t_{\mathrm{disp}}| \approx 6 \times 10^{-62}$--$1 \times 10^{-60}$\,s. That is $56$--$57$ orders of magnitude below a $10^{-4}$\,s timing benchmark for third-generation detectors. A delay of $10^{-4}$\,s at $z = 3$ would require $\ell_\ast \approx 3 \times 10^{-7}$\,m, i.e.\ an energy scale $\hbar c/\ell_\ast \approx 0.7$\,eV instead of $M_P$. Such a scale would have to be confronted with existing dispersion constraints from LIGO--Virgo--KAGRA and would not follow from the Planck-scale model.

\subsection{An Ansatz for the Running of the Tensor Tilt}
Let $n_t = \dif\ln\mathcal{P}_t/\dif\ln k$ be the tilt of the primordial tensor power spectrum $\mathcal{P}_t(k)$, and let $\alpha_t \coloneqq \dif n_t/\dif\ln k$ be its running. Let $p$ be a physical momentum and $M_\ast$ the energy scale of the effect. The two-point Pad\'e form $d_s(p) = 2 + 2/(1 + (p/M_\ast)^2)$ has the same limits $2$ and $4$ as \eqref{eq:analytical_ds_formula}, but it is an interpolation, not that closed form. Adopting the ansatz $\alpha_t = \frac{1}{2}(d_s(p) - 4)$, which does not follow from the diffusion model, gives
\begin{equation}
    \alpha_t = \frac{\dif n_t}{\dif\ln k} = -\frac{(p/M_\ast)^2}{1 + (p/M_\ast)^2}.
\end{equation}
We take $M_\ast = M_P$. A local quantum-gravity effect with scale $M_P$ must compare $M_P$ with the \emph{physical} momentum $p$ of a tensor mode when the mode is generated, which is at horizon exit, $p = k/a = H_{\mathrm{inf}}$; the comoving wavenumber measured today is not an invariant scale. The bound $r < 0.036$ on the tensor-to-scalar ratio \cite{bicepkeck2021constraints} and the scalar amplitude $A_s \approx 2.1 \times 10^{-9}$ \cite{planck2018parameters} give, through $2H_{\mathrm{inf}}^2/(\pi^2 M_{\mathrm{red}}^2) = r A_s$, the bound $H_{\mathrm{inf}} \le 4.7 \times 10^{13}$\,GeV. Hence $|\alpha_t| \approx (H_{\mathrm{inf}}/M_P)^2 \le 1.5 \times 10^{-11}$ (or $3.7 \times 10^{-10}$ with $M_{\mathrm{red}}$), nearly the same for all CMB multipoles because $H_{\mathrm{inf}}$ varies little while they exit the horizon. The planned experiments LiteBIRD and CMB-S4 target the tensor amplitude at the level $r \sim 10^{-3}$, far from any sensitivity to a running of this size. The ansatz therefore offers no observational target.

\subsection{Quantum Simulators}
Programmable quantum simulators (Rydberg arrays, superconducting qubits) can measure quantum Fisher metrics of tensor-network states and out-of-time-order correlators. Such measurements test properties of specific quantum many-body states. They do not test the conjectures of this chapter unless a concrete model is exhibited for which the conjectured relations (for instance Conjecture~\ref{thm:horizon_morse_saturation}) make a quantitative prediction.

\subsection{Atom Interferometry}
Atom interferometry does not test a ``Jordan-loop protection'' effect. The restriction to Jordan loops is a conjecture about a kinematical sector (Conjecture~\ref{conj:jordan_sector}), not a consequence of causality, and for strontium atoms $\ell_P/\lambda_{\mathrm{dB}} = m_{\mathrm{Sr}} v\, \ell_P/(2\pi\hbar) \approx 3.5 \times 10^{-27}\,(v/1\,\mathrm{m\,s^{-1}})$, which stays below $2 \times 10^{-25}$ for the launch velocities ($v \lesssim 45$\,m\,s$^{-1}$) of long-baseline instruments.

% =========================================================================
\section{Theoretical Correspondences}
\label{sec:master_table}
% =========================================================================

Table~\ref{tab:grand_unification} establishes the theoretical dictionary uniting all four domains of the theory.

\begin{table}[htbp]
\centering
\scriptsize
\renewcommand{\arraystretch}{1.35}
\resizebox{\textwidth}{!}{%
\begin{tabular}{@{}llll@{}}
\toprule
\textbf{Simplicial Calculus (\emph{Part II})} & \textbf{Minimax Geometry (\emph{Part III})} & \textbf{Tensor Networks (\emph{Part I})} & \textbf{Spacetime Physics (GR/QFT)} \\
\midrule
Continuous Pascal Simplex $\Delta_m(x)$ & Space-like Cauchy Foliation $\Sigma_t$ & Tensor-Train Variety $\mathcal{M}_{\mathbf{r}}^{\mathrm{TT}}$ & Space-time Slicing $(M^4, g_{\mu\nu})$ \\
Fractional Laplacian $\Delta_{\Delta_m}^\alpha$ & Minimax Extrinsic Curvature $\|\II\|_{L^\infty}$ & Projected TT Gradient Flow & Regulated Wheeler--DeWitt $\mathcal{H}=0$ \\
Running Spectral Dim $d_s(\alpha) \to 2$ & Bounded Extrinsic Curvature $K_{ij}K^{ij} \le 3\kappa^2$ & Entanglement Renormalization & Planckian UV Renormalizability \\
Lie $A_{m-1}$ Cartan Metric & Causal Jordan Embeddings & Continuous Wilson Loops $\Tr(\mathcal{P}e^{\oint A})$ & Ashtekar Connection $A_a^i \in \mathfrak{su}(2)$ \\
Beta-Kernel Convolutions & Israel Thin-Shell Junctions ($C^{0,1}$ metric) & Fubini--Study Metric $g^{\mathrm{FS}}$ & Emergent AdS Metric $\dif s^2 = \frac{L^2}{z^2}(\dif z^2 + \dif x^2)$ \\
Polytope Boundary Recurrences & Boundary $\partial \Omega$ Obstacles & Entanglement Cut Minimization & Ryu--Takayanagi Minimal Area $\frac{\Area}{4 G}$ \\
Combinatorial Multinomial Entropy & Medial Axis Variational Geometry & Quantum Relative Entropy $S(\rho \| \sigma)$ & Einstein Equations $G_{\mu\nu} = 8\pi G T_{\mu\nu}$ \\
Sierpinski Fractal Measure & Self-Intersection Ban & Graphon Ricci Neckpinch & Topological Foam Condensation \\
Dispersive Simplicial Waves & Proper Four-Acceleration $a^\mu$ & Kac--Rice Tensor Morse Complexity & Black Hole Microstate Scrambling \\
\bottomrule
\end{tabular}%
}
\caption{Theoretical correspondences between Simplicial Fractional Calculus, Minimax Relativistic Foliations, Multilinear Tensor Networks, and Spacetime Gravitation. Entries in the same row are heuristic analogies, not isomorphisms; the only links proved in this chapter are those of Propositions~\ref{thm:running_spectral_dimension}, \ref{thm:cartan_metric_emergence}, \ref{thm:neckpinch_condensation} (conditional on its curvature hypothesis) and Theorem~\ref{thm:minimax_hamiltonian_regularization}; Theorem~\ref{thm:grand_einstein_emergence} is imported from the literature.}
\label{tab:grand_unification}
\end{table}

% =========================================================================
\section{Conclusion and the Road to Planckian Quantum Cosmology}
\label{sec:conclusion}
% =========================================================================

We have separated what the three programs of this treatise establish about quantum gravity from what they suggest.
\begin{enumerate}
    \item \textbf{Established:} closed-form spectral dimensions for homogeneous and two-scale diffusion symbols, including an exact interpolation $d_s: 2 \to 4$ for the symbol $k^2 + \ell_P^2 k^4$; sharp pointwise bounds on ${}^{(3)}R$ and $K_{ij}K^{ij}$ on slices of bounded extrinsic curvature; the identification of the barycentric Hessian of the continuous multinomial with the $A_{m-1}$ root metric; and the exponential decay of graphon entries on which the Ollivier curvature stays below $-c < 0$, a conditional statement whose hypothesis depends on the metric used to define the curvature (Remark after Proposition~\ref{thm:neckpinch_condensation}) and has not been verified for a specific kernel.
    \item \textbf{Imported from the literature:} under the standard holographic hypotheses, the linearized and second-order entanglement derivation of the Einstein equations (Theorem~\ref{thm:grand_einstein_emergence}).
    \item \textbf{Conjectured:} the invariance and anomaly-free closure of a Jordan-loop sector of loop quantum gravity; mean curvature flow as a search procedure for Ryu--Takayanagi surfaces (not every initial cut flows to the minimizer); the condensation of graphon Ricci flows to smooth four-dimensional geometries; and the identification of Kac--Rice complexity with horizon entropy.
    \item \textbf{Observationally:} Planck-suppressed signatures of the kind considered here are many orders of magnitude below experimental reach (Section~\ref{sec:experimental}).
\end{enumerate}
The framework does not, at present, constitute a UV-complete or anomaly-free theory of quantum gravity. Proving the conjectures above is the natural programme for future work.

% Shared declarations block, \input by every chapter before its bibliography.
% Keep in sync with the "Declarations" section of master_book_unified_quantum_gravity.tex.
\section*{Declarations}

\noindent\textbf{Funding.} This research was financed in part by the Coordena\c{c}\~ao de Aperfei\c{c}oamento de Pessoal de N\'ivel Superior -- Brasil (CAPES) -- Finance Code 001.

\medskip\noindent\textbf{Competing interests.} The author declares no competing interests.

\medskip\noindent\textbf{Use of generative AI tools.} Large language model assistants (Anthropic Claude; Google Gemini, used through the Antigravity agent environment) were used extensively in preparing this work: drafting and editing text and \LaTeX{} source; proposing, expanding and revising derivations and proofs under the author's direction; writing the Python verification scripts and the \textsc{Lean 4} files; searching and checking the literature and references; and adversarial auditing of proofs, numerical claims and code. These audits found errors, some of them originally introduced with AI assistance. The errors and their corrections are recorded in the repository file \texttt{WORKPLAN.md}. AI tools are not authors. The author chose the research questions and the overall framework, reviewed the material, and is responsible for the content, including any remaining errors.

\medskip\noindent\textbf{Verification status.} Results labelled Theorem, Proposition or Lemma are proved in the text or cited as classical; statements labelled Conjecture are not proved. The numerical scripts compare formulas of the text with independent computations and are checks, not proofs. The \textsc{Lean 4} modules in \texttt{formal\_proofs\_book/Book} are a naming skeleton and do not verify the mathematics of this chapter.

\medskip\noindent\textbf{Code and data availability.} Sources, numerical scripts and the audit log are available at
\begin{center}\url{https://github.com/reinaldomsilvafilho-netizen/quantum-gravity-lean4}.\end{center}
The monograph is archived on Zenodo, \href{https://doi.org/10.5281/zenodo.22290043}{doi:10.5281/zenodo.22290043}, under the CC~BY~4.0 license.


\begin{thebibliography}{99}

\bibitem{ambjorn2005spectral}
J.~Ambj{\o}rn, J.~Jurkiewicz, and R.~Loll,
\emph{The spectral dimension of the universe is scale dependent},
Physical Review Letters, \textbf{95} (2005), no.~17, 171301,
\doi{10.1103/PhysRevLett.95.171301}.

\bibitem{ashtekar2004background}
A.~Ashtekar and J.~Lewandowski,
\emph{Background independent quantum gravity: a status report},
Classical and Quantum Gravity, \textbf{21} (2004), no.~15, R53,
\doi{10.1088/0264-9381/21/15/R01}.

\bibitem{faulkner2014gravitation}
T.~Faulkner, M.~Guica, T.~Hartman, R.~C. Myers, and M.~Van~Raamsdonk,
\emph{Gravitation from entanglement in holographic CFTs},
Journal of High Energy Physics, \textbf{2014} (2014), no.~3, 51,
\doi{10.1007/JHEP03(2014)051}.

\bibitem{faulkner2017nonlinear}
T.~Faulkner, F.~M. Haehl, E.~Hijano, O.~Parrikar, C.~Rabideau, and M.~Van~Raamsdonk,
\emph{Nonlinear gravity from entanglement in conformal field theories},
Journal of High Energy Physics, \textbf{2017} (2017), no.~8, 57,
\doi{10.1007/JHEP08(2017)057}.

\bibitem{horava2009spectral}
P.~Ho\v{r}ava,
\emph{Spectral dimension of the universe in quantum gravity at a Lifshitz point},
Physical Review Letters, \textbf{102} (2009), 161301,
\doi{10.1103/PhysRevLett.102.161301}.

\bibitem{sotiriou2011spectral}
T.~P. Sotiriou, M.~Visser, and S.~Weinfurtner,
\emph{Spectral dimension as a probe of the ultraviolet continuum regime of causal dynamical triangulations},
Physical Review Letters, \textbf{107} (2011), 131303,
\doi{10.1103/PhysRevLett.107.131303}.

\bibitem{sotiriou2011dispersion}
T.~P. Sotiriou, M.~Visser, and S.~Weinfurtner,
\emph{From dispersion relations to spectral dimension---and back again},
Physical Review D, \textbf{84} (2011), 104018,
\doi{10.1103/PhysRevD.84.104018}.

\bibitem{carlip2017dimension}
S.~Carlip,
\emph{Dimension and dimensional reduction in quantum gravity},
Classical and Quantum Gravity, \textbf{34} (2017), 193001,
\doi{10.1088/1361-6382/aa8535}.

\bibitem{maldacena2016remarks}
J.~Maldacena and D.~Stanford,
\emph{Remarks on the Sachdev--Ye--Kitaev model},
Physical Review D, \textbf{94} (2016), 106002,
\doi{10.1103/PhysRevD.94.106002}.

\bibitem{auffinger2013random}
A.~Auffinger, G.~Ben Arous, and J.~{\v{C}}ern{\'y},
\emph{Random matrices and complexity of spin glasses},
Communications on Pure and Applied Mathematics, \textbf{66} (2013), no.~2, 165--201,
\doi{10.1002/cpa.21422}.

\bibitem{lashkari2016canonical}
N.~Lashkari and M.~Van~Raamsdonk,
\emph{Canonical energy is quantum Fisher information},
Journal of High Energy Physics, \textbf{2016} (2016), no.~4, 153,
\doi{10.1007/JHEP04(2016)153}.

\bibitem{maldacena2016bound}
J.~Maldacena, S.~H. Shenker, and D.~Stanford,
\emph{A bound on chaos},
Journal of High Energy Physics, \textbf{2016} (2016), no.~8, 106,
\doi{10.1007/JHEP08(2016)106}.

\bibitem{ryu2006holographic}
S.~Ryu and T.~Takayanagi,
\emph{Holographic derivation of entanglement entropy from the anti-de Sitter space/conformal field theory correspondence},
Physical Review Letters, \textbf{96} (2006), no.~18, 181602,
\doi{10.1103/PhysRevLett.96.181602}.

\bibitem{silvafilho2026flows} R.~M. Silva-Filho, \emph{Geometric Flows, Partial Differential Equations, and Variational Dynamics on Matrix and Tensor Manifolds}, Chapter~2 of \emph{Geometry, Tensors, and Quantum Gravity}, Universidade Federal de Lavras, 2026.

\bibitem{silvafilho2026functional} R.~M. Silva-Filho, \emph{Functional Realizations of Matrices and Hypertensors: Emergent Invariants and Geometric Measures}, Chapter~1 of \emph{Geometry, Tensors, and Quantum Gravity}, Universidade Federal de Lavras, 2026.

\bibitem{silvafilho2026noneuclidean} R.~M. Silva-Filho, \emph{Minimax Extrinsic Curvature in Non-Euclidean Geometries and General Relativity: Variational Theory, Space Forms, and ADM Spacetime Slicings}, Chapter~8 of \emph{Geometry, Tensors, and Quantum Gravity}, Universidade Federal de Lavras, 2026.

\bibitem{silvafilho2026pascal} R.~M. Silva-Filho, \emph{Analytic Continuation of Pascal's Simplex: Continuous Multinomial Integrals, Polytope Boundary Recurrences, and Simplicial Fractional Calculus}, Chapter~3 of \emph{Geometry, Tensors, and Quantum Gravity}, Universidade Federal de Lavras, 2026.

\bibitem{silvafilho2026spacetime} R.~M. Silva-Filho, \emph{Emergent Spacetime and Quantum Geometry: Tensor Networks, Holonomies, and Entanglement (Established Results and Conjectures)}, Chapter~11 of \emph{Geometry, Tensors, and Quantum Gravity}, Universidade Federal de Lavras, 2026.

\bibitem{swingle2012entanglement}
B.~Swingle,
\emph{Entanglement renormalization and holography},
Physical Review D, \textbf{86} (2012), no.~6, 065007,
\doi{10.1103/PhysRevD.86.065007}.

\bibitem{verstraete2010continuous}
F.~Verstraete and J.~I. Cirac,
\emph{Continuous matrix product states for quantum fields},
Physical Review Letters, \textbf{104} (2010), no.~19, 190405,
\doi{10.1103/PhysRevLett.104.190405}.

\bibitem{koide1981}
Y.~Koide,
\emph{Fermion-boson two-body model of quarks and leptons and Cabibbo mixing},
Lettere al Nuovo Cimento, \textbf{34} (1982), 201--205,
\href{https://doi.org/10.1007/BF02817096}{DOI: 10.1007/BF02817096}.

\bibitem{thiemann1998qsd}
T.~Thiemann,
\emph{Quantum spin dynamics (QSD)},
Classical and Quantum Gravity, \textbf{15} (1998), 839--873,
\doi{10.1088/0264-9381/15/4/011}.

\bibitem{swingle2014universality}
B.~Swingle and M.~Van~Raamsdonk,
\emph{Universality of gravity from entanglement},
preprint (2014), arXiv:1405.2933,
\doi{10.48550/arXiv.1405.2933}.

\bibitem{deharo2001holographic}
S.~de~Haro, K.~Skenderis, and S.~N. Solodukhin,
\emph{Holographic reconstruction of spacetime and renormalization in the AdS/CFT correspondence},
Communications in Mathematical Physics, \textbf{217} (2001), 595--622,
\doi{10.1007/s002200100381}.

\bibitem{chamseddine2007gravity}
A.~H. Chamseddine, A.~Connes, and M.~Marcolli,
\emph{Gravity and the standard model with neutrino mixing},
Advances in Theoretical and Mathematical Physics, \textbf{11} (2007), 991--1089,
\doi{10.4310/ATMP.2007.v11.n6.a3}.

\bibitem{jacob2008lorentz}
U.~Jacob and T.~Piran,
\emph{Lorentz-violation-induced arrival delays of cosmological particles},
Journal of Cosmology and Astroparticle Physics, \textbf{2008} (2008), no.~1, 031,
\doi{10.1088/1475-7516/2008/01/031}.

\bibitem{mirshekari2012constraining}
S.~Mirshekari, N.~Yunes, and C.~M. Will,
\emph{Constraining Lorentz-violating, modified dispersion relations with gravitational waves},
Physical Review D, \textbf{85} (2012), 024041,
\doi{10.1103/PhysRevD.85.024041}.

\bibitem{bicepkeck2021constraints}
P.~A.~R. Ade et~al. (BICEP/Keck Collaboration),
\emph{Improved constraints on primordial gravitational waves using Planck, WMAP, and BICEP/Keck observations through the 2018 observing season},
Physical Review Letters, \textbf{127} (2021), 151301,
\doi{10.1103/PhysRevLett.127.151301}.

\bibitem{planck2018parameters}
N.~Aghanim et~al. (Planck Collaboration),
\emph{Planck 2018 results. VI. Cosmological parameters},
Astronomy \& Astrophysics, \textbf{641} (2020), A6,
\doi{10.1051/0004-6361/201833910}.

\end{thebibliography}

\end{document}
